\chapter{随机振动初步}
\intro{随机振动研究激励和响应具有随机性的系统振动问题。本章介绍随机过程的基本概念、统计描述方法、频域分析及其在单自由度线弹性系统中的应用。}

\section{随机过程基础}
\subsection{确定性振动与随机振动的对比}
\textbf{确定性振动}：未来任意时刻的响应可由运动方程和初始条件精确预测，例如简谐激励下单自由度系统的稳态响应。

\textbf{随机振动}：激励本身是不确定的时间函数，无法精确预测任一特定时间历程。只能从统计意义上描述振动行为的宏观特征。典型实例包括：
\begin{itemize}
    \item 地震地面运动加速度
    \item 路面不平度对车辆的激励
    \item 风荷载湍流分量
    \item 海洋波浪随机载荷
    \item 喷气噪声引起的结构振动
\end{itemize}

\subsection{随机过程的定义}
\begin{definition}[随机过程]
设 $(\Omega, \mathcal{F}, P)$ 为概率空间，$T \subseteq \R$ 为时间指标集。随机过程 $\{X(t), t \in T\}$ 是从 $\Omega \times T$ 到 $\R$ 的二元映射，满足对固定 $t$，$X(t, \cdot)$ 是随机变量;对固定 $\omega$，$X(\cdot, \omega)$ 是实数函数（样本函数或实现）。
\end{definition}

\begin{itemize}
    \item \textbf{样本函数（Sample Function）}：固定 $\omega$ 得到的时间函数 $x(t)$，对应一次试验的观测结果
    \item \textbf{系集（Ensemble）}：全部样本函数的集合，代表随机过程的总体
    \item \textbf{系集平均（Ensemble Average）}：对固定 $t$，对全部样本函数取期望
\end{itemize}

\subsection{平稳随机过程}
\begin{definition}[严平稳（Strict-sense Stationary）]
若随机过程的有限维联合分布对任意时移 $\tau$ 保持不变：
\[
F_{X}(x_1,\dots,x_n; t_1+\tau, \dots, t_n+\tau) = F_{X}(x_1,\dots,x_n; t_1, \dots, t_n)
\]
称该过程为严平稳过程。
\end{definition}

工程中更常用的是\textbf{宽平稳（Weak-sense Stationary）}，仅要求：
\begin{enumerate}
    \item 均值不随时间变化：$\mu_X(t) = \mu_X = \text{const}$
    \item 自相关函数仅与时间差有关：$R_X(t_1, t_2) = R_X(\tau)$，$\tau = t_2 - t_1$
\end{enumerate}

\subsection{各态历经性（Ergodicity）}
\begin{definition}[各态历经性]
若随机过程的时间平均（沿单一样本函数）等于系集平均（沿全部样本函数）：
\[
\lim_{T\to\infty} \frac{1}{2T}\int_{-T}^{T} x(t)\,dt = \mu_X = E[X(t)]
\]
\[
\lim_{T\to\infty} \frac{1}{2T}\int_{-T}^{T} x(t)x(t+\tau)\,dt = R_X(\tau) = E[X(t)X(t+\tau)]
\]
则称该过程具有\textbf{各态历经性}。各态历经过程必为平稳过程。
\end{definition}

各态历经性的重要性在于：工程师通常仅观测到一次实现（如一次地震记录），各态历经假设允许我们使用单样本的时间平均来推断其统计特征。

\section{统计描述}
\subsection{均值与均方值}
\begin{definition}[均值函数]
\[
\mu_X(t) = E[X(t)] = \int_{-\infty}^{\infty} x \, f_X(x, t)\,dx
\]
其中 $f_X(x,t)$ 为 $t$ 时刻的一阶概率密度函数。对于平稳过程 $\mu_X = \text{const}$。
\end{definition}

\begin{definition}[均方值]
\[
\psi_X^{2}(t) = E[X^{2}(t)] = \int_{-\infty}^{\infty} x^{2} \, f_X(x, t)\,dx
\]
均方值的平方根 $\psi_X$ 称为\textbf{均方根值（RMS）}。
\end{definition}

\subsection{方差与标准差}
\begin{definition}[方差函数]
\[
\sigma_X^{2}(t) = E[(X(t) - \mu_X(t))^{2}] = E[X^{2}(t)] - \mu_X^{2}(t) = \psi_X^{2}(t) - \mu_X^{2}(t)
\]
$\sigma_X$ 为标准差（Standard Deviation），反映随机变量相对均值的离散程度。
\end{definition}

\begin{myformula}
对于零均值过程（$\mu_X = 0$），方差等于均方值：
\[
\sigma_X^{2} = \psi_X^{2} = E[X^{2}]
\]
这在分析振动信号时极为常见——通常先对信号去趋势处理（消除均值），剩余部分即为零均值过程。
\end{myformula}

\subsection{自相关函数}
\begin{definition}[自相关函数]
\[
R_X(t_1, t_2) = E[X(t_1) X(t_2)] = \int_{-\infty}^{\infty}\int_{-\infty}^{\infty} x_1 x_2\, f_X(x_1, x_2; t_1, t_2)\, dx_1 dx_2
\]
对于宽平稳过程，$R_X(\tau) = E[X(t)X(t+\tau)]$。
\end{definition}

自相关函数的性质：
\begin{enumerate}
    \item $R_X(0) = E[X^{2}] = \psi_X^{2} \ge 0$（零时差给出均方值）
    \item $R_X(\tau) = R_X(-\tau)$（偶函数）
    \item $|R_X(\tau)| \le R_X(0)$（零时差取最大值）
    \item $\lim_{\tau \to \infty} R_X(\tau) = \mu_X^{2}$（若过程不包含周期成分）
    \item 若 $X(t)$ 含周期成分，则 $R_X(\tau)$ 也含同频周期成分
\end{enumerate}

\begin{figure}[H]
\centering
\begin{tikzpicture}[scale=1.0]
    \draw[-Stealth] (-4,0) -- (4,0) node[right] {$\tau$};
    \draw[-Stealth] (0,-0.5) -- (0,3.5) node[above] {$R_X(\tau)$};

    \draw[blue,very thick] plot[domain=-3.8:3.8,samples=200] (\x,{3*exp(-abs(\x))});
    \draw[red,thick,dashed] plot[domain=-3.8:3.8,samples=200] (\x,{3*exp(-abs(\x))*cos(300*\x)});
    \draw[green!50!black,thick] plot[domain=-3.8:3.8,samples=200] (\x,{3*exp(-1.5*abs(\x))});

    \node[blue] at (2.5,1.2) {宽带（慢衰减）};
    \node[green!50!black] at (2.5,2.2) {宽带（快衰减）};
    \node[red] at (2.5,0.3) {窄带（含周期）};

    \draw[gray,dashed] (-4,{3*exp(-4)}) -- (0,{3*exp(-4)});
    \node[gray] at (-3.2,0.6) {$\mu_X^{2}$};
    \node at (0.3,3.0) {$\psi_X^{2}$};
    \filldraw (0,3) circle (1.5pt);
\end{tikzpicture}
\caption{自相关函数示意。蓝色：指数衰减（宽带随机）;红色：指数衰减加余弦调制（窄带，含优势频率）;绿色：更快衰减（更大带宽）。}
\label{fig:ch11_autocorr}
\end{figure}

\subsection{互相关函数}
\begin{definition}[互相关函数]
对两个联合宽平稳过程 $X(t)$ 和 $Y(t)$：
\[
R_{XY}(\tau) = E[X(t)Y(t+\tau)]
\]
\[
R_{YX}(\tau) = E[Y(t)X(t+\tau)] = R_{XY}(-\tau)
\]
互相关函数测量两个过程的时间延迟关联，当两个过程独立时 $R_{XY}(\tau) = \mu_X \mu_Y$。
\end{definition}

\section{频域描述}
\subsection{功率谱密度（Power Spectral Density, PSD）}
\begin{definition}[功率谱密度]
对宽平稳随机过程 $X(t)$，功率谱密度函数（双侧谱）定义为自相关函数的 Fourier 变换：
\[
S_X(\omega) = \int_{-\infty}^{\infty} R_X(\tau) e^{-i\omega\tau} d\tau
\]
逆变换由 Wiener-Khinchin 定理给出：
\[
R_X(\tau) = \frac{1}{2\pi}\int_{-\infty}^{\infty} S_X(\omega) e^{i\omega\tau} d\omega
\]
\end{definition}

\begin{myformula}
零时差的自相关函数等于 PSD 的积分（总均方值）：
\[
\psi_X^{2} = R_X(0) = \frac{1}{2\pi}\int_{-\infty}^{\infty} S_X(\omega)\,d\omega
\]
建议注意单侧谱 $W_X(\omega) = 2 S_X(\omega)$（$\omega \ge 0$ 时），此时：
\[
\psi_X^{2} = \frac{1}{2\pi}\int_{0}^{\infty} W_X(\omega)\,d\omega = \int_{0}^{\infty} W_X(f)\,df
\]
其中 $W_X(f) = 2\pi W_X(\omega)/\pi$ 为工程常用频率尺度单侧 PSD（Hz）。
\end{myformula}

PSD 的物理意义：$S_X(\omega_0)\Delta\omega$ 表示频带 $[\omega_0, \omega_0 + \Delta\omega]$ 内信号功率的贡献。对于位移过程，PSD 单位为 $[\text{m}^{2}/\text{Hz}]$（面积/频率）;对于加速度，单位为 $[(\text{m}/\text{s}^{2})^{2}/\text{Hz}]$ 或常用 dB 尺度。

\subsection{白噪声（White Noise）}
\begin{definition}[白噪声]
白噪声是理想化的随机过程，其 PSD 在所有频率上为常数：
\[
S_X(\omega) = S_0 = \text{const}, \quad -\infty < \omega < \infty
\]
相应的自相关函数为：
\[
R_X(\tau) = S_0\,\delta(\tau)
\]
其中 $\delta(\tau)$ 为 Dirac $\delta$ 函数。白噪声任意两不同时刻取值完全无关，均方值 $\psi_X^{2} \to \infty$（物理不可实现）。
\end{definition}

物理近似：若随机过程 PSD 在远大于系统分析带宽的频率范围内保持平直，则可近似为白噪声。热噪声（Johnson-Nyquist 噪声）在 $\sim 10^{12}\,\text{Hz}$ 以下可视为白噪声。

\subsection{窄带与宽带随机过程}
\begin{itemize}
    \item \textbf{宽带过程}：PSD 分布在较宽的频率范围内，对应的自相关函数衰减迅速
    \item \textbf{窄带过程}：能量集中在以中心频率 $\omega_0$ 为中心的窄频带内，对应的自相关函数呈现衰减的余弦振荡：
    \[
    R_X(\tau) = \alpha e^{-\beta|\tau|}\cos(\omega_0 \tau)
    \]
    典型的窄带过程为线性系统在宽带激励下的共振响应
\end{itemize}

\begin{figure}[H]
\centering
\begin{tikzpicture}[scale=1.0]
    \begin{scope}
        \draw[-Stealth] (0,0) -- (5,0) node[right] {$\omega$};
        \draw[-Stealth] (0,0) -- (0,3.5) node[above] {$S_X(\omega)$};
        \draw[blue,very thick] (0,2.5) -- (4.5,2.5);
        \node[blue] at (2.5,3.0) {白噪声};
        \draw[gray,dashed] (0,2.5) node[left] {$S_0$};
    \end{scope}

    \begin{scope}[xshift=6cm]
        \draw[-Stealth] (0,0) -- (5,0) node[right] {$\omega$};
        \draw[-Stealth] (0,0) -- (0,3.5) node[above] {$S_X(\omega)$};
        \draw[red,very thick] plot[domain=0:5,samples=100] (\x,{2.8*exp(-0.3*(\x-2.5)^2)});
        \node[red] at (2.5,3.0) {窄带噪声};
    \end{scope}

    \begin{scope}[xshift=12cm]
        \draw[-Stealth] (0,0) -- (5,0) node[right] {$\omega$};
        \draw[-Stealth] (0,0) -- (0,3.5) node[above] {$S_X(\omega)$};
        \draw[green!50!black,very thick] plot[domain=0.1:4.8,samples=80]
            (\x,{2.0*(1 - 0.6/(1 + 0.5*(\x-2.2)^2))});
        \node[green!50!black] at (2.5,3.0) {宽带噪声};
    \end{scope}
\end{tikzpicture}
\caption{白噪声、窄带噪声与宽带噪声的功率谱密度对比。}
\label{fig:ch11_psd_types}
\end{figure}

\section{线性系统对随机激励的响应}
\subsection{频率响应函数回顾}
线性时不变单自由度系统在力 $f(t)$ 激励下的运动方程：
\[
m\ddot{x} + c\dot{x} + kx = f(t)
\]
系统的\textbf{频率响应函数（FRF）} $H(\omega)$ 定义为单位幅值简谐激励下的稳态响应：
\[
H(\omega) = \frac{1}{k - m\omega^{2} + i c\omega} = \frac{1/k}{1 - (\omega/\omega_n)^{2} + i\,2\zeta(\omega/\omega_n)}
\]
其中 $\omega_n = \sqrt{k/m}$, $\zeta = c/(2m\omega_n)$。

\subsection{输入输出 PSD 关系}
\begin{theorem}[随机激励下的传递关系]
若 $X(t)$ 为宽平稳随机激励，系统为线性时不变系统，则稳态响应 $Y(t)$ 也为宽平稳随机过程，其功率谱密度满足：
\[
S_Y(\omega) = |H(\omega)|^{2} S_X(\omega)
\]
互功率谱密度：
\[
S_{XY}(\omega) = H(\omega) S_X(\omega)
\]
\end{theorem}

\begin{myformula}
响应均方值为 PSD 在全频域的积分：
\[
E[Y^{2}] = \frac{1}{2\pi}\int_{-\infty}^{\infty} |H(\omega)|^{2} S_X(\omega)\,d\omega
\]
这是随机振动分析的核心公式——系统的频率选择特性 $|H(\omega)|^{2}$ 对输入谱施加加权，总均方响应为各频率分量贡献的叠加。
\end{myformula}

\subsection{单自由度系统对白噪声激励的响应}
设输入为理想白噪声 $S_X(\omega) = S_0$，则响应均方值为：
\[
E[Y^{2}] = \frac{S_0}{2\pi}\int_{-\infty}^{\infty} |H(\omega)|^{2}\,d\omega
\]

\begin{myformula}
利用有理函数积分公式（对二次分母的积分）：
\[
\int_{-\infty}^{\infty} \frac{d\omega}{(\omega_n^{2} - \omega^{2})^{2} + (2\zeta\omega_n\omega)^{2}} = \frac{\pi}{2\zeta\omega_n^{3}}
\]
得单自由度白噪声基座激励下的位移均方响应：
\[
E[Y^{2}] = \frac{\pi S_0}{2\zeta\omega_n^{3} m^{2}} = \frac{\pi S_0}{c\,k}
\]
可见增大阻尼 $c$ 和刚度 $k$ 均可有效压缩均方响应。
\end{myformula}

\begin{remark}
白噪声激励下，单自由度系统的位移均方响应与阻尼成反比（$E[Y^{2}] \propto 1/\zeta$）。这给出了一种随机振动减振策略：在共振频带附近引入附加阻尼的效果最为显著，因为 $|H(\omega)|^{2}$ 在该区域取最大值。
\end{remark}

\subsection{单侧谱与工程量纲}
在工程振动中，常用加速度 PSD（常以 $g^{2}/\text{Hz}$ 为量纲）表征激励。对位移系统的白噪声基座加速度激励 $\ddot{x}_g$：
\[
m\ddot{x} + c\dot{x} + kx = -m\ddot{x}_g
\]
设 $S_{\ddot{x}_g}(\omega) = S_a$（常数，单位 $[\text{m}^{2}/\text{s}^{4}/\text{Hz}]$），则：
\[
S_x(\omega) = \frac{S_a}{(\omega_n^{2} - \omega^{2})^{2} + (2\zeta\omega_n\omega)^{2}}
\]
响应加速度 PSD：
\[
S_{\ddot{x}}(\omega) = \omega^{4} S_x(\omega) = \frac{\omega^{4} S_a}{(\omega_n^{2} - \omega^{2})^{2} + (2\zeta\omega_n\omega)^{2}}
\]

\begin{figure}[H]
\centering
\begin{tikzpicture}[scale=1.0]
    \draw[-Stealth] (0,0) -- (6,0) node[right] {$\omega/\omega_n$};
    \draw[-Stealth] (0,0) -- (0,4) node[above] {$|H(\omega)|^{2}$};

    \draw[green!50!black,very thick,dashed] (0,2.5) -- (5.5,2.5);
    \node[green!50!black] at (4,2.9) {$S_X$ = const};

    \draw[blue,very thick] plot[smooth,domain=0.05:5.5,samples=200]
        (\x,{2.5/((1-\x*\x)^2 + 4*0.05*\x*\x)});
    \node[blue] at (2.0,3.6) {$|H(\omega)|^{2} S_X(\omega)$};

    \draw[gray,dashed] (1,0) -- (1,4);
    \node at (1.0,4.2) {$\omega_n$};
\end{tikzpicture}
\caption{白噪声输入（绿色虚线，PSD 为常数）经单自由度系统传递后，输出 PSD（蓝色曲线）在固有频率 $\omega_n$ 附近出现尖峰——系统充当窄带滤波器。}
\label{fig:ch11_input_output_PSD}
\end{figure}

\section{疲劳与首次穿越问题}
\subsection{峰值分布与穿越率}
随机振动可导致结构的疲劳失效。对于零均值的窄带 Gaussian 过程 $X(t)$，其峰值概率密度由 Rayleigh 分布描述：
\[
f_P(p) = \frac{p}{\sigma_X^{2}}\exp\left(-\frac{p^{2}}{2\sigma_X^{2}}\right), \quad p \ge 0
\]
其中 $\sigma_X^{2}$ 为过程的方差。

\begin{definition}[穿越率（Crossing Rate）]
单位时间内过程向上穿越水平 $x = \alpha$ 的期望次数为（Rice 公式）：
\[
\nu_{\alpha}^{+} = \nu_0 \exp\left(-\frac{\alpha^{2}}{2\sigma_X^{2}}\right)
\]
其中 $\nu_0$ 为零穿越率：
\[
\nu_0 = \frac{1}{2\pi}\frac{\sigma_{\dot{X}}}{\sigma_X} = \frac{1}{2\pi}\sqrt{\frac{\int_0^{\infty} \omega^{2} S_X(\omega)\,d\omega}{\int_0^{\infty} S_X(\omega)\,d\omega}}
\]
$\nu_0$ 近似等于过程的表观频率（apparent frequency）。
\end{definition}

\subsection{Miner 线性累积损伤理论}
构件在变幅应力循环下的疲劳累积用 Miner 法则：
\[
D = \sum_i \frac{n_i}{N_i}
\]
其中 $n_i$ 为应力水平 $S_i$ 下的实际循环次数，$N_i$ 为该水平下的疲劳寿命（由 S-N 曲线查得）。$D = 1$ 时发生疲劳失效。

对于随机振动，窄带过程在时间 $T$ 内的总应力循环次数近似为 $\nu_0 T$，其幅值分布由 Rayleigh 分布给出，累积损伤为：
\[
D \approx \nu_0 T \int_0^{\infty} \frac{f_P(s)}{N(s)}\,ds
\]
代入 Basquin 形式的 S-N 曲线 $N(s) = C s^{-b}$ 后可得到闭合的疲劳寿命预测公式。

\subsection{首次穿越问题（First-Passage Problem）}
首次穿越问题研究随机过程首次超过阈值（安全边界）的概率和时间分布。设阈值 $x = b$，定义可靠性函数：
\[
R(t) = P\{\max_{0 \le s \le t} X(s) < b\}
\]
对于 Gaussian 窄带过程且 $b/\sigma_X \gg 1$，首次穿越时间近似服从指数分布（Poisson 穿越假设）：
\[
R(t) \approx \exp\left(-\nu_b^{+} t\right)
\]
失效率（hazard rate）由穿越率 $\nu_b^{+}$ 直接给出。

\section{Python 代码实现}
\subsection{白噪声生成与单自由度随机响应}
\begin{mycode}{白噪声激励下单自由度系统的响应模拟}
\begin{lstlisting}[language=Python]
import numpy as np
import matplotlib.pyplot as plt
from scipy.signal import welch, lsim, TransferFunction

np.random.seed(42)
fs = 200.0
T = 100.0
N = int(T * fs)
t = np.linspace(0, T, N)

wn = 2.0 * np.pi * 2.0
zeta = 0.05
m = 1.0
k = m * wn**2
c = 2 * zeta * m * wn

num = [1.0]
den = [m, c, k]
sys = TransferFunction(num, den)

S0 = 0.5
dt = 1.0 / fs
white_noise = np.sqrt(S0 * fs) * np.random.randn(N)

t_resp, x, _ = lsim(sys, U=white_noise, T=t)

f, Pxx = welch(white_noise, fs, nperseg=4096)
f, Pyy = welch(x, fs, nperseg=4096)

fig, axes = plt.subplots(2, 2, figsize=(12, 8))
axes[0,0].plot(t[-2000:], x[-2000:], lw=0.4)
axes[0,0].set_title('位移响应时程（尾部）')
axes[0,0].set_xlabel('$t$ [s]'); axes[0,0].set_ylabel('$x(t)$')

axes[0,1].hist(x, bins=60, density=True, alpha=0.7, color='blue')
axes[0,1].set_title('位移概率密度')
axes[0,1].set_xlabel('$x$'); axes[0,1].set_ylabel('PDF')

axes[1,0].loglog(f, Pxx, lw=0.6, label='输入')
axes[1,0].loglog(f, Pyy, lw=0.8, label='输出')
axes[1,0].legend(); axes[1,0].set_title('PSD 对比 (Welch)')
axes[1,0].set_xlabel('$f$ [Hz]'); axes[1,0].set_ylabel('$S(f)$')

f_theory = np.linspace(0.01, 5, 500)
w = 2.0 * np.pi * f_theory
H2 = np.abs(1.0 / (k - m*w**2 + 1j*c*w))**2
axes[1,1].plot(f_theory, S0 * H2, lw=1.5, label='理论 $|H|^2 S_0$')
axes[1,1].plot(f, Pyy, lw=0.6, alpha=0.7, label='模拟 Welch')
axes[1,1].legend(); axes[1,1].set_title('输出 PSD 理论与模拟对比')
axes[1,1].set_xlabel('$f$ [Hz]'); axes[1,1].set_ylabel('$S_x(f)$')
axes[1,1].set_xlim(0, 5)
plt.tight_layout(); plt.show()
\end{lstlisting}
\end{mycode}

\subsection{PSD 估计（Welch 方法）与穿越率计算}
\begin{mycode}{随机响应信号的穿越率估计}
\begin{lstlisting}[language=Python]
def crossing_rate(x, dt, level=0.0):
    upward_crossings = ((x[:-1] < level) & (x[1:] >= level)).sum()
    return upward_crossings / (len(x) * dt)

def spectral_moment(Sx, f, k):
    df = f[1] - f[0]
    return np.trapz((2*np.pi*f)**k * Sx, f)

f, Pxx = welch(x, fs, nperseg=4096)
m0 = spectral_moment(Pxx, f, 0)
m2 = spectral_moment(Pxx, f, 2)
nu0_theory = np.sqrt(m2 / m0) / (2 * np.pi)
nu0_sim = crossing_rate(x, dt)
print(f"理论零穿越率: {nu0_theory:.3f} Hz")
print(f"模拟零穿越率: {nu0_sim:.3f} Hz")

b_levels = np.linspace(0.5, 3.0, 6) * np.sqrt(m0)
nu_b = np.array([crossing_rate(x, dt, b) for b in b_levels])
for bv, nv in zip(b_levels / np.sqrt(m0), nu_b):
    print(f"  b/sigma = {bv:.1f}, nu_b^+ = {nv:.4f} Hz")
\end{lstlisting}
\end{mycode}

\section{习题与思考}
\begin{enumerate}
    \item 证明宽平稳过程自相关函数 $R_X(\tau)$ 的偶对称性 $R_X(\tau) = R_X(-\tau)$，并证明 $R_X(0) \ge |R_X(\tau)|$。
    \item 设单自由度系统 $m=1$, $\omega_n=10\,\text{rad/s}$, $\zeta=0.05$，在理想白噪声力激励 $S_F(\omega)=1$ 下运行。计算位移和速度的均方值，并分别写出位移和速度功率谱密度的表达式。
    \item 推导窄带 Gaussian 过程的峰值 Rayleigh 分布，并结合 Rice 公式给出阈值穿越率的表达式。
    \item 用 Python 生成一组平稳 Gaussian 白噪声、窄带噪声和宽带噪声，绘制各自的时程、直方图、自相关函数和 PSD 估计图，对比验证理论性质。
    \item 解释：为何白噪声在物理上不可实现？通常采用什么模型代替？说明滤波白噪声的基本思想与实现方法。
    \item 理解各态历经假设在工程随机振动信号处理中的关键作用——为何常仅用一次记录就能估算系集统计量？
    \item 基于 Welche PSD 估计，编程计算并对比不同参数（分段长度、窗函数）对 PSD 估计精度和分辨率的影响。
\end{enumerate}
